\subsection{约化定理}

我们保持第 2 节至第 7 节的记号与假设。

引理 6. 设 R 是交换环，\(\mathfrak{p}_i\) （ \(1 \leqslant i \leqslant n\) ）是 R 的互异素理想。

\begin{enumerate}
\def\labelenumi{(\roman{enumi})}
\item
  对 \(1 \leqslant i \leqslant n\) ，设 \(\mathbf{H}_i\) 是 \(\mathbf{R} / \mathfrak{p}_i\) 的满足下列条件的子集：不存在元素 \(\alpha_i \in \mathbf{R} / \mathfrak{p}_i\) ，使 \(\alpha_i + \mathbf{H}_i\) 含有某个 \(\neq 0\) 理想，即 \(\mathbf{R} / \mathfrak{p}_i\) 的一个理想。则存在 \(a \in \mathbf{R}\) ，使对 \(1 \leqslant i \leqslant n\) ，\(a\) 在 \(\mathbf{R} / \mathfrak{p}_i\) 中的典范像不属于 \(\mathbf{H}_i\) 。
\item
  若 \(\operatorname{Card}(\mathbf{H}_i) < \operatorname{Card}(\mathbf{R} / \mathfrak{p}_i)\) ，则诸 \(\mathbf{H}_i\) 满足 (i) 中的条件。
\item
  我们对 \(n\) 作归纳，\(n = 0\) 的情形是平凡的。于是设 \(n \geqslant 1\) 。我们可以对指标 \(i\) 作一个置换，从而可设 \(p_1\) 在诸 \(p_i\) 中极小，于是对 \(2 \leqslant i \leqslant n\) ，存在 \(c_i \in p_i\) 使 \(c_i \notin p_1\) 。由归纳假设，存在 \(b \in R\) ，使 \(b\) 在 \(R / p_i\) 中的典范像不属于 \(H_i\) （对 \(2 \leqslant i \leqslant n\) ）。对一切 \(x \in R\) ，我们写 \(a_{x}=b+xc_{2}c_{3}\ldots c_{n};\) 由于 \(c_{i}\in\mathfrak{p}_{i},\) 显然 \(a_{x}\equiv b\pmod{\mathfrak{p}_{i}}\) （对 \(2\leqslant i\leqslant n\) ）。因此只需证明存在 \(x\in R\) ，使 \(a_{x}\) 在 \(R/\mathfrak{p}_{1}\) 中的典范像不属于 \(H_{1}\) 。现在，当 x 跑遍 R 时，诸 \(a_{x}\) 在 \(R/\mathfrak{p}_{1}\) 中的典范像所成的集正是 \(\beta+c\) ，其中 \(\beta\) 是 b 的典范像，c 是 \(R/\mathfrak{p}_{1}\) 中由 \(c_{2}c_{3}\ldots c_{n}\) 的典范像生成的理想；鉴于诸 \(c_{i},c\neq0\) ，因为 \(R/\mathfrak{p}_{1}\) 是整环，而关于 \(H_{1}\) 的假设蕴含存在一个解决问题的 x。
\item
  由于 \(R/p_{i}\) 是整环，每个 \(\neq0\) 理想（\(R/p_{i}\) 中的）的基数都等于 \(R/p_{i}\) 的基数，且任一理想经 \(R/p_{i}\) 的元素所作的平移亦然，由此得结论。
\end{enumerate}

定理 6. 设 M 是无扭的有限生成 A-模。存在 M 的自由子模 L，使 M/L 同构于 A 的一个理想。

我们用 \(n\) 表示 \(M\) 的秩（即 \(V = M \otimes_{A} K\) 在 \(K\) 上的秩），并把 \(M\) 视为 \(V\) 关于 \(A\) 的格。于是对一切 \(p \in P(A)\) ，\(M_p\) 是 \(V\) 关于 \(A_p\) 的格（第 1 节，例 6），且由于 \(A_p\) 是主理想整环，\(M_p\) 是自由 \(A_p\) -模，秩为 \(n\) 。我们将写：

\[
\mathbf {M} (\mathfrak {p}) = \mathbf {M} _ {\mathfrak {p}} / \mathfrak {p} \mathbf {M} _ {\mathfrak {p}}.
\]

我们用 \(k(\mathfrak{p})\) 表示 \(\mathbf{A} / \mathfrak{p}\) 的分式域（同构于 \(\mathbf{A}_{\mathfrak{p}}\) 的剩余域）；于是 \(\mathbf{M}(\mathfrak{p}) = \mathbf{M} \otimes_{\mathbf{A}} k(\mathfrak{p})\) 是秩为 \(n\) 的向量空间（在 \(k(\mathfrak{p})\) 上）。对一切 \(x \in \mathbf{M}\) ，我们用 \(x(\mathfrak{p})\) 表示 \(x\) 在 \(\mathbf{M}(\mathfrak{p})\) 中的典范像。

引理 7. 设 \(x_{i}\) （ \(1 \leqslant i \leqslant m\) ）是 \(\mathbf{M}\) 中（在 \(\mathbf{A}\) 上或在 \(\mathbf{K}\) 上，二者等价）线性无关的元素，\(\mathbf{L}\) 是 \(\mathbf{M}\) 的由诸 \(x_{i}\) 生成的子 A-模。则对几乎所有 \(\mathfrak{p} \in \mathbf{P}\) ，\(x_{i}(\mathfrak{p}) \in \mathbf{M}(\mathfrak{p})\) 在 \(k(\mathfrak{p})\) 上线性无关；它们在 \(k(\mathfrak{p})\) 上对一切 \(\mathfrak{p} \in \mathbf{P}\) 线性无关，当且仅当 \(\mathbf{M}/\mathbf{L}\) 无扭。

设 \(x_{m+1}, \ldots, x_n\) 是 M 中与 \(x_1, \ldots, x_m\) 一起构成 V 的基的元素，N 是 M 的由诸 \(x_i (1 \leqslant i \leqslant n)\) 生成的自由子 A-模。由第 3 节的定理 3，\(N_p = M_p\) 对几乎所有 \(p \in P\) 成立；由于 \(x_1(p), \ldots, x_n(p)\) 构成 \(N(p)\) 在 \(k(p)\) 上的基，这就证明了第一个断言。

若 \(\mathbf{M} / \mathbf{L}\) 无扭，则 \((\mathbf{M} / \mathbf{L})_{\mathfrak{p}} = \mathbf{M}_{\mathfrak{p}} / \mathbf{L}_{\mathfrak{p}}\) 对一切 \(\mathfrak{p} \in \mathbb{P}\) 也无扭（第 1 节，例 6），且由于 \(\mathbf{A}_{\mathfrak{p}}\) 是主理想整环，\(\mathbf{M}_{\mathfrak{p}} / \mathbf{L}_{\mathfrak{p}}\) 是自由的。因此 \(\mathbf{M}_{\mathfrak{p}}\) 是 \(\mathbf{L}_{\mathfrak{p}}\) 与一个自由 \(\mathbf{A}_{\mathfrak{p}}\) -模 E （秩为 \(n - m\) ）的直和；于是 \(\mathbf{M}(\mathfrak{p})\) 是 \(\mathbf{L}(\mathfrak{p})\) 与 \(k(\mathfrak{p})\) -向量空间 \(\mathrm{E} / \mathfrak{p}\mathrm{E}\) （秩为 \(n - m\) ）的直和；因此 \(\mathbf{L}(\mathfrak{p})\) 的秩为 \(m\) ，而它由诸 \(x_{i}(\mathfrak{p})\) （ \(1 \leqslant i \leqslant m\) ）生成，故这些元素线性无关。

反之，设诸 \(x_{i}(\mathfrak{p})\) （ \(1 \leqslant i \leqslant m\) ）在 \(k(\mathfrak{p})\) 上对一切 \(\mathfrak{p} \in \mathbb{P}\) 线性无关。则 \(\mathbf{L}_{\mathfrak{p}}\) 是 \(\mathbf{M}_{\mathfrak{p}}\) 的直因子（对一切 \(\mathfrak{p}\) ）（第 II 章，§3，第 2 节，命题 5 的推论 1），因此 \(\mathbf{M}_{\mathfrak{p}} / \mathbf{L}_{\mathfrak{p}} = (\mathbf{M} / \mathbf{L})_{\mathfrak{p}}\) 对一切 \(p \in P\) 无扭。由第 IV 章，§1，第 2 节，命题 5 的推论，我们断定 \(\mathrm{P} \cap \mathrm{Ass}(\mathrm{M}/\mathrm{L}) = \varnothing\)。但由于 L 是自反的，由第 2 节命题 7 (i)，能属于 \(\mathrm{Ass}(\mathrm{M}/\mathrm{L})\) 的唯一素理想是理想 (0)；故 M/L 无扭。

引理 8. 设 M 的秩 n \(\geqslant2\) ；则存在 M 的元素 \(x \neq 0\) ，使 M/Ax 无扭。

设 \(y \neq 0\) 是 M 的一个元素。由引理 7，\(p \in P\) 中使 \(y(p) = 0\) 的元素所成的集 Y 是有限的。若 Y = \(\varnothing\) ，则由引理 7（应用于仅由单独一个元素 y 组成的序列 \((x_i)\) ）可知 M/Ay 无扭。因此设 Y ≠ \(\varnothing\) ，并写 S = \(\bigcap_{p \in Y} (A - p)\) ；我们知道（第 4 节，引理 2），\(S^{-1}A\) 是半局部主理想整环，其极大理想是诸 \(pS^{-1}A\) （p ∈ Y），相应的局部环是诸 \(A_{p}\) 。因此

\[
\mathrm{S} ^ {- 1} \mathrm{A} / \mathfrak {p} \mathrm{S} ^ {- 1} \mathrm{A} = k (\mathfrak {p}),
\]

由此得

\[
\begin{array}{c} \mathrm{S} ^ {- 1} \mathrm{M} / \mathfrak {p} \mathrm{S} ^ {- 1} \mathrm{M} = (\mathrm{M} / \mathfrak {p} \mathrm{M}) \otimes_ {\mathrm{A}} \mathrm{S} ^ {- 1} \mathrm{A} = \mathrm{M} \otimes_ {\mathrm{A}} ((\mathrm{A} / \mathfrak {p}) \otimes_ {\mathrm{A}} \mathrm{S} ^ {- 1} \mathrm{A}) = \\ \mathrm{M} \otimes_ {\mathrm{A}} k (\mathfrak {p}) = \mathrm{M} (\mathfrak {p}) \end{array}
\]

对一切 \(\mathfrak{p} \in \mathbf{Y}\) 。由第 II 章，§ 1，第 2 节，命题 6，存在元素 \(z / s \in \mathrm{S}^{-1}\mathbf{M}\) （ \(z \in \mathbf{M}, s \in \mathrm{S}\) ），它在诸 \(\mathbf{M}(\mathfrak{p})\) （ \(\mathfrak{p} \in \mathbf{Y}\) ）中的典范像都 \(\neq 0\) 。于是由 \(S\) 的定义，\(z(\mathfrak{p}) \neq 0\) 对一切 \(\mathfrak{p} \in Y\) 成立。我们还可以进一步设 \(y\) 与 \(z\) 在 \(K\) 上线性无关。因为在相反的情形，考虑元素 \(t \in M\) ，它与 \(y\) 线性无关（由于 \(n \geqslant 2\) ，这样的元素存在）；另一方面取一个元素 \(a \neq 0\) ，它属于 \(\bigcap_{\mathfrak{p} \in Y} \mathfrak{p}\) （由于 \(A\) 是整环，它不化为 0），并写 \(z' = z + at\) ：显然 \(y\) 与 \(z'\) 在 \(K\) 上线性无关，且 \(z'(\mathfrak{p}) = z(\mathfrak{p}) \neq 0\) 对一切 \(\mathfrak{p} \in Y\) 成立。

于是设 \(y\) 与 \(z\) 在 \(\mathbf{K}\) 上线性无关，令 \(\mathbf{Z}\) 为 \(\mathfrak{p} \in \mathbf{P} - \mathbf{Y}\) 中使 \(y(\mathfrak{p})\) 与 \(z(\mathfrak{p})\) 在 \(k(\mathfrak{p})\) 上线性无关者所成的集；由引理 7，这个集是有限的。对一切 \(\mathfrak{p} \in \mathbf{Z}\) ，我们可以写 \(z(\mathfrak{p}) = \lambda(\mathfrak{p})y(\mathfrak{p})\) ，其中 \(\lambda(\mathfrak{p}) \in k(\mathfrak{p})\) 。现在，\(\operatorname{Card}(\mathrm{A}/\mathfrak{p}) \geqslant 2\) 对一切 \(\mathfrak{p} \in \mathbf{P}\) 成立；因此由引理 6，存在 \(b \in \mathbf{A}\) ，使对一切 \(\mathfrak{p} \in \mathbf{Z}\) ，\(b\) 在 \(\mathrm{A}/\mathfrak{p}\) 中的典范像异于 \(\lambda(\mathfrak{p})\) 。我们现在证明元素 \(x = z - by\) 解决问题；只需（利用取 \(m = 1\) 的引理 7）验证 \(x(\mathfrak{p}) \neq 0\) 对一切 \(\mathfrak{p} \in \mathbf{P}\) 成立。现在：

--- 若 \(\mathfrak{p} \in \mathbf{Y}\) ，则按构造 \(x(\mathfrak{p}) \neq 0\) ；

--- 若 \(p \in Z\) ，则 \(x(p) = \mu_{p} y(p)\) ，其中由 b 的选取 \(\mu_{p} \neq 0\) ，故 \(x(p) \neq 0\) ，因为 \(y(p) \neq 0\) ；

- 若 \(\mathfrak{p} \in \mathrm{P} - (\mathrm{Y} \cup \mathrm{Z}), y(\mathfrak{p})\) 与 \(z(\mathfrak{p})\) 线性无关，故 \(x(\mathfrak{p}) \neq 0\) 。

这些引理既已建立，我们转向定理 6 的证明。我们对 n 作归纳，\(n \leqslant 1\) 的情形是平凡的，因为此时 M 本身同构于 A 的一个理想。设 \(n \geqslant 2\) ；由引理 8，存在 M 的秩 1 自由子模 \(L_{0}\) ，使 \(M/L_{0}\) 无扭；于是 \(M/L_{0}\) 的秩为 n - 1。由归纳假设，存在自由子模 \(L_{1}\) ，它含于 \(M/L_{0}\) ，使 \((\mathrm{M}/\mathrm{L}_{0})/\mathrm{L}_{1}\) 同构于 A 的一个理想。设 L 是 \(L_{1}\) 在 M 中的逆像；\(L/L_{0}\) 同构于 \(L_{1}\) ，且由于 \(L_{1}\) 自由，L 同构于 \(L_{0} \oplus L_{1}\) （《代数》第 II 章，§ 1，第 11 节，命题 21），从而是自由的；由于 M/L 同构于 \((\mathrm{M}/\mathrm{L}_{0})/\mathrm{L}_{1}\) ，定理得证。

注. 若 M 是自反的，M/L 未必自反（习题 9）。

